
\documentclass{article}
\usepackage{colortbl}
\usepackage{makecell}
\usepackage{multirow}
\usepackage{supertabular}

\begin{document}

\newcounter{utterance}

\centering \large Interaction Transcript for game `cladder', experiment `full\_v1.5\_default', episode 2927 with llama{-}3.1{-}8b{-}instruct.
\vspace{24pt}

{ \footnotesize  \setcounter{utterance}{1}
\setlength{\tabcolsep}{0pt}
\begin{supertabular}{c@{$\;$}|p{.15\linewidth}@{}p{.15\linewidth}p{.15\linewidth}p{.15\linewidth}p{.15\linewidth}p{.15\linewidth}}
    \# & \multicolumn{2}{c}{Player} && \multicolumn{2}{c}{Game Master} \\
    \hline

    \theutterance \stepcounter{utterance}  
    & & & \multicolumn{4}{p{0.6\linewidth}}{
        \cellcolor[rgb]{0.9,0.9,0.9}{
            \makecell[{{p{\linewidth}}}]{
                \texttt{\tiny{[P1$\langle$GM]}}
                \texttt{You are an expert in causal inference. The following question is not a typical commonsense query, but rather a meticulously designed question created by a professor specializing in causal inference, intended to assess the students' mastery of the course content.} \\
\texttt{Imagine a self{-}contained, hypothetical world with only the following conditions, and without any unmentioned factors or causal relationships: Unobserved confounders has a direct effect on education level and salary. Proximity to a college has a direct effect on education level. Education level has a direct effect on salary. Unobserved confounders is unobserved. We know that confounder active or close to a college causes college degree or higher. confounder active and college degree or higher causes high salary. We observed the person lives close to a college and confounder active. Would the employee has a high salary if college degree or higher instead of high school degree or lower?} \\
\texttt{Start your answer with 'Yes' or 'No', followed by additional reasoning or evidence to support your explanation.} \\
            }
        }
    }
     \\ \\

    \theutterance \stepcounter{utterance}  
    & \multicolumn{4}{p{0.6\linewidth}}{
        \cellcolor[rgb]{0.561,0.737,0.561}{
            \makecell[{{p{\linewidth}}}]{
                \texttt{\tiny{[P1$\rangle$GM]}}
                \texttt{No} \\
\\ 
\texttt{In this hypothetical world, we have the following causal relationships:} \\
\\ 
\texttt{1. Unobserved confounders (U) affects education level (E) and salary (S).} \\
\texttt{2. Proximity to a college (P) affects education level (E).} \\
\texttt{3. Education level (E) affects salary (S).} \\
\\ 
\texttt{Given the observed information: } \\
\\ 
\texttt{{-} The person lives close to a college (P).} \\
\texttt{{-} The confounder is active (U).} \\
\\ 
\texttt{We can infer that the person has a college degree or higher (E) due to the condition that confounder active or close to a college causes college degree or higher.} \\
\\ 
\texttt{However, to determine the effect of having a college degree or higher on salary, we need to consider the potential effect of the confounder (U) on salary. Since the confounder is active, it directly affects salary (S). } \\
\\ 
\texttt{Given the condition that confounder active and college degree or higher causes high salary, it might seem like having a college degree or higher would lead to a high salary. However, this is a conditional effect that assumes the confounder is active. } \\
\\ 
\texttt{In the absence of the confounder's effect, we cannot conclude that having a college degree or higher would lead to a high salary. The observed effect of college degree or higher on salary is confounded by the active confounder (U). } \\
\\ 
\texttt{Therefore, without additional information or assumptions about the effect of the confounder (U) on salary in the absence of a college degree or higher, we cannot conclude that the employee would have a high salary if they had a college degree or higher instead of a high school degree or lower.} \\
            }
        }
    }
    & & \\ \\

    \theutterance \stepcounter{utterance}  
    & & & \multicolumn{2}{p{0.3\linewidth}}{
        \cellcolor[rgb]{0.9,0.9,0.9}{
            \makecell[{{p{\linewidth}}}]{
                \texttt{\tiny{[GM$|$GM]}}
                \texttt{No} \\
            }
        }
    }
    & & \\ \\

    \theutterance \stepcounter{utterance}  
    & & & \multicolumn{2}{p{0.3\linewidth}}{
        \cellcolor[rgb]{0.9,0.9,0.9}{
            \makecell[{{p{\linewidth}}}]{
                \texttt{\tiny{[GM$|$GM]}}
                \texttt{yes} \\
            }
        }
    }
    & & \\ \\

    \theutterance \stepcounter{utterance}  
    & & & \multicolumn{2}{p{0.3\linewidth}}{
        \cellcolor[rgb]{0.9,0.9,0.9}{
            \makecell[{{p{\linewidth}}}]{
                \texttt{\tiny{[GM$|$GM]}}
                \texttt{game\_result = LOSE} \\
            }
        }
    }
    & & \\ \\

\end{supertabular}
}

\end{document}
